\begin{thebibliography}{99}
\bibitem{source1} Karl Abrahamson. Time-space tradeoffs for algebraic problems on general sequential machines. Journal of Computer and System Sciences, 43(2):269–289, 1991. doi:10.1016/0022-0000(91) 90014-V.
\bibitem{source2} N. Alon, J. Bruck, J. Naor, M. Naor, and R.M. Roth. Construction of asymptotically good low- rate error-correcting codes through pseudo-random graphs. IEEE Transactions on Information Theory, 38(2):509–516, 1992. doi:10.1109/18.119713.
\bibitem{source3} Laasya Bangalore, Rishabh Bhadauria, Carmit Hazay, and Muthuramakrishnan Venkitas- ubramaniam. On black-box constructions of time and space efficient sublinear arguments from symmetric-key primitives. In Theory of Cryptography: 20th International Conference, TCC 2022, Chicago, IL, USA, November 7–10, 2022, Proceedings, Part I, pages 417–446. Springer-Verlag, 2022. doi:10.1007/978-3-031-22318-1\_15.
\bibitem{source4} L.M.J. Bazzi and S.K. Mitter. Endcoding complexity versus minimum distance. IEEE Transactions on Information Theory, 51(6):2103–2112, 2005. doi:10.1109/TIT.2005.847727.
\bibitem{source5} P. Beame. A general sequential time-space tradeoff for finding unique elements. In Proceedings of the Twenty-First Annual ACM Symposium on Theory of Computing, STOC ’89, pages 197–203, New York, NY, USA, 1989. Association for Computing Machinery. doi:10.1145/73007.73026.
\bibitem{source6} Paul Beame and Niels Kornerup. Cumulative memory lower bounds for randomized and quantum computation. In Kousha Etessami, Uriel Feige, and Gabriele Puppis, editors, 50th International Colloquium on Automata, Languages, and Programming (ICALP 2023), volume
\bibitem{source261} of Leibniz International Proceedings in Informatics (LIPIcs), pages 17:1–17:20, Dagstuhl, Germany, 2023. Schloss Dagstuhl – Leibniz-Zentrum für Informatik. doi:10.4230/LIPIcs. ICALP.2023.17.
\bibitem{source7} Paul Beame, Michael Saks, Xiaodong Sun, and Erik Vee. Time-space trade-off lower bounds for randomized computation of decision problems. J. ACM, 50(2):154–195, 2003. doi: 10.1145/636865.636867.
\bibitem{source8} Eli Ben-Sasson, Alessandro Chiesa, Daniel Genkin, and Eran Tromer. On the concrete efficiency of probabilistically-checkable proofs. In Proceedings of the Forty-Fifth Annual ACM Symposium on Theory of Computing, STOC ’13, pages 585–594. Association for Computing Machinery, 2013. doi:10.1145/2488608.2488681.
\bibitem{source9} Charles H. Bennett. Time/space trade-offs for reversible computation. SIAM J. Comput., 18(4):766–776, 1989. doi:10.1137/0218053.
\bibitem{source10} Nir Bitansky and Alessandro Chiesa. Succinct arguments from multi-prover interactive proofs and their efficiency benefits. In Proceedings of the 32nd Annual Cryptology Conference on Advances in Cryptology — CRYPTO 2012 - Volume 7417, pages 255–272, Berlin, Heidelberg, 2012. Springer-Verlag. doi:10.1007/978-3-642-32009-5\_16.
\bibitem{source11} A. Borodin and S. Cook. A time-space tradeoff for sorting on a general sequential model of computation. SIAM J. Comput., 11(2):287–297, 1982. doi:10.1137/0211022.
\bibitem{source12} M. Capalbo, O. Reingold, S. Vadhan, and A. Wigderson. Randomness conductors and constant- degree lossless expanders. In Proceedings 17th IEEE Annual Conference on Computational Complexity, pages 8–8, 2002. doi:10.1109/CCC.2002.1004327.
\bibitem{source13} Eshan Chattopadhyay, Vipul Goyal, and Xin Li. Non-malleable extractors and codes, with their many tampered extensions. In Proceedings of the Forty-Eighth Annual ACM Symposium on Theory of Computing, STOC ’16, pages 285–298. Association for Computing Machinery, 2016. doi:10.1145/2897518.2897547.
\bibitem{source14} M. Cheraghchi, F. Didier, and A. Shokrollahi. Invertible extractors and wiretap protocols. IEEE Trans. Inf. Theor., 58(2):1254–1274, 2012. doi:10.1109/TIT.2011.2170660.
\bibitem{source15} Mahdi Cheraghchi and Venkatesan Guruswami. Non-malleable coding against bit-wise and split-state tampering. J. Cryptol., 30(1):191–241, 2017. doi:10.1007/s00145-015-9219-z.
\bibitem{source16} Joshua Cook and Dana Moshkovitz. Time and space efficient deterministic decoders. Elec- tronic Colloquium on Computational Complexity, 2024. URL: https://eccc.weizmann.ac. il/report/2024/110/.
\bibitem{source17} Yotam Dikstein, Irit Dinur, and Shiri Sivan. The linear time encoding scheme fails to encode, 2023. arXiv:2312.16125.
\bibitem{source18} Irit Dinur, Shai Evra, Ron Livne, Alexander Lubotzky, and Shahar Mozes. Locally testable codes with constant rate, distance, and locality. In Proceedings of the 54th Annual ACM SIGACT Symposium on Theory of Computing, STOC 2022, pages 357–374. Association for Computing Machinery, 2022. doi:10.1145/3519935.3520024.
\bibitem{source19} Dariush Divsalar, Hui Jin, and Robert J. McEliece. Coding theorems for ’turbo-like’ codes. In Proceedings 36th Annual Allerton Conference on Communication, Control, and Computing, pages 201–210, 1998. URL: https://api.semanticscholar.org/CorpusID:1045655.
\bibitem{source20} R. Gallager. Low-density parity-check codes. IRE Transactions on Information Theory, 8(1):21–28, 1962. doi:10.1109/TIT.1962.1057683.
\bibitem{source21} Shafi Goldwasser, Yael Tauman Kalai, and Guy N. Rothblum. Delegating computation: Interactive proofs for muggles. J. ACM, 62(4), September 2015. doi:10.1145/2699436.
\bibitem{source22} André Gronemeier. A note on the decoding complexity of error-correcting codes. Inf. Process. Lett., 100(3):116–119, 2006. doi:10.1016/j.ipl.2006.06.006.
\bibitem{source23} V. Guruswami and P. Indyk. Linear-time encodable/decodable codes with near-optimal rate. IEEE Transactions on Information Theory, 51(10):3393–3400, 2005. doi:10.1109/TIT.2005. 855587.
\bibitem{source24} Venkatesan Guruswami and Widad Machmouchi. Explicit interleavers for a repeat accumulate accumulate (raa) code construction. In 2008 IEEE International Symposium on Information Theory, pages 1968–1972, 2008. doi:10.1109/ISIT.2008.4595333.
\bibitem{source25} Venkatesan Guruswami, Christopher Umans, and Salil Vadhan. Unbalanced expanders and randomness extractors from parvaresh-vardy codes. In Twenty-Second Annual IEEE Conference on Computational Complexity (CCC’07), pages 96–108, 2007. doi:10.1109/CCC.2007.38.
\bibitem{source26} Anna Gál, Kristoffer Arnsfelt Hansen, Michal Koucký, Pavel Pudlák, and Emanuele Viola. Tight bounds on computing error-correcting codes by bounded-depth circuits with arbitrary gates. IEEE Transactions on Information Theory, 59(10):6611–6627, 2013. doi:10.1109/TIT. 2013.2270275.
\bibitem{source27} R. W. Hamming. Error detecting and error correcting codes. The Bell System Technical Journal, 29(2):147–160, 1950. doi:10.1002/j.1538-7305.1950.tb00463.x.
\bibitem{source28} Tzvika Hartman and Ran Raz. On the distribution of the number of roots of polynomials and explicit weak designs. Random Struct. Algorithms, 23(3):235–263, October 2003. doi: 10.1002/rsa.10095.
\bibitem{source29} F.C. Hennie. One-tape, off-line turing machine computations. Information and Control, 8(6):553–578, 1965. doi:10.1016/S0019-9958(65)90399-2.
\bibitem{source30} Justin Holmgren and Ron Rothblum. Delegating computations with (almost) minimal time and space overhead. In 2018 IEEE 59th Annual Symposium on Foundations of Computer Science (FOCS), pages 124–135, 2018. doi:10.1109/FOCS.2018.00021.
\bibitem{source31} R. Impagliazzo, L. A. Levin, and M. Luby. Pseudo-random generation from one-way functions. In Proceedings of the Twenty-First Annual ACM Symposium on Theory of Computing, STOC ’89, pages 12–24. Association for Computing Machinery, 1989. doi:10.1145/73007.73009.
\bibitem{source32} S. Jukna. A nondeterministic space-time tradeoff for linear codes. Information Processing Letters, 109(5):286–289, 2009. doi:10.1016/j.ipl.2008.11.001.
\bibitem{source33} Itay Kalev and Amnon Ta-Shma. Unbalanced Expanders from Multiplicity Codes. In Amit Chakrabarti and Chaitanya Swamy, editors, Approximation, Randomization, and Combin- atorial Optimization. Algorithms and Techniques (APPROX/RANDOM 2022), volume 245 of Leibniz International Proceedings in Informatics (LIPIcs), pages 12:1–12:14, Dagstuhl, Germany, 2022. Schloss Dagstuhl – Leibniz-Zentrum für Informatik. doi:10.4230/LIPIcs. APPROX/RANDOM.2022.12.
\bibitem{source34} Swastik Kopparty, Or Meir, Noga Ron-Zewi, and Shubhangi Saraf. High-rate locally-correctable and locally-testable codes with sub-polynomial query complexity. In Proceedings of the Forty-Eighth Annual ACM Symposium on Theory of Computing, STOC ’16, pages 202–215. Association for Computing Machinery, 2016. doi:10.1145/2897518.2897523.
\bibitem{source35} Swastik Kopparty, Shubhangi Saraf, and Sergey Yekhanin. High-rate codes with sublinear-time decoding. J. ACM, 61(5), 2014. doi:10.1145/2629416.
\bibitem{source36} S. Rao Kosaraju. Real-time simulation of concatenable double-ended queues by double-ended queues (preliminary version). In Proceedings of the Eleventh Annual ACM Symposium on Theory of Computing, STOC ’79, pages 346–351. Association for Computing Machinery, 1979. doi:10.1145/800135.804427.
\bibitem{source37} Xin Li. Improved non-malleable extractors, non-malleable codes and independent source extractors. In Proceedings of the 49th Annual ACM SIGACT Symposium on Theory of Computing, STOC 2017, pages 1144–1156. Association for Computing Machinery, 2017. doi: 10.1145/3055399.3055486.
\bibitem{source38} D. J. C. Mackay and Radford M. Neal. Near shannon limit performance of low density parity check codes. Electronics Letters, 33:457–458, 1996. URL: https://api.semanticscholar. org/CorpusID:122801915.
\bibitem{source39} Noam Nisan and David Zuckerman. Randomness is linear in space. J. Comput. Syst. Sci., 52(1):43–52, 1996. doi:10.1006/jcss.1996.0004.
\bibitem{source40} Wolfgang J. Paul, Joel I. Seiferas, and Janos Simon. An information-theoretic approach to time bounds for on-line computation (preliminary version). In Proceedings of the Twelfth Annual ACM Symposium on Theory of Computing, STOC ’80, pages 357–367. Association for Computing Machinery, 1980. doi:10.1145/800141.804685.
\bibitem{source41} Ran Raz, Omer Reingold, and Salil Vadhan. Extracting all the randomness and reducing the error in trevisan’s extractors. In Proceedings of the Thirty-First Annual ACM Symposium on Theory of Computing, STOC ’99, pages 149–158. Association for Computing Machinery, 1999. doi:10.1145/301250.301292.
\bibitem{source42} I.S. Reed. A brief history of the development of error correcting codes. Computers \& Mathematics with Applications, 39(11):89–93, 2000. doi:10.1016/S0898-1221(00)00112-7.
\bibitem{source43} O. Reingold, R. Shaltiel, and A. Wigderson. Extracting randomness via repeated condensing. In Proceedings 41st Annual Symposium on Foundations of Computer Science, pages 22–31, 2000. doi:10.1109/SFCS.2000.892008.
\bibitem{source44} Omer Reingold, Guy N. Rothblum, and Ron D. Rothblum. Constant-round interactive proofs for delegating computation. In Proceedings of the Forty-Eighth Annual ACM Symposium on Theory of Computing, STOC ’16, pages 49–62. Association for Computing Machinery, 2016. doi:10.1145/2897518.2897652.
\bibitem{source45} Nandakishore Santhi and Alexander Vardy. Minimum distance of codes and their branching program complexity. In 2006 IEEE International Symposium on Information Theory, pages 1490–1494, 2006. doi:10.1109/ISIT.2006.262116.
\bibitem{source46} Walter J. Savitch and Paul M. B. Vitányi. Linear time simulation of multihead turing machines with head-to-head jumps. In Proceedings of the Fourth Colloquium on Automata, Languages and Programming, pages 453–464, Berlin, Heidelberg, 1977. Springer-Verlag.
\bibitem{source47} M. Sipser and D.A. Spielman. Expander codes. IEEE Transactions on Information Theory, 42(6):1710–1722, 1996. doi:10.1109/18.556667.
\bibitem{source48} D.A. Spielman. Linear-time encodable and decodable error-correcting codes. IEEE Transactions on Information Theory, 42(6):1723–1731, 1996. doi:10.1109/18.556668.
\bibitem{source49} Daniel A. Spielman. Linear-time encodable and decodable error-correcting codes. In Proceedings of the Twenty-Seventh Annual ACM Symposium on Theory of Computing, STOC ’95, pages 388–397. Association for Computing Machinery, 1995. doi:10.1145/225058.225165.
\bibitem{source50} Amnon Ta-Shma, Christopher Umans, and David Zuckerman. Loss-less condensers, unbalanced expanders, and extractors. In Proceedings of the Thirty-Third Annual ACM Symposium on Theory of Computing, STOC ’01, pages 143–152. Association for Computing Machinery, 2001. doi:10.1145/380752.380790.
\bibitem{source51} Luca Trevisan. Construction of extractors using pseudo-random generators (extended abstract). In Proceedings of the Thirty-First Annual ACM Symposium on Theory of Computing, STOC ’99, pages 141–148. Association for Computing Machinery, 1999. doi:10.1145/301250.301289.
\bibitem{source52} Yaacov Yesha. Time-space tradeoffs for matrix multiplication and the discrete fourier transform on any general sequential random-access computer. Journal of Computer and System Sciences, 29(2):183–197, 1984. doi:10.1016/0022-0000(84)90029-1.
\end{thebibliography}
